% !PRINTED 604
\weberpage{604}
%<<CONT
corresponding basis form $\alpha_1 t_1 + \alpha_2 t_2$ that is associated with $\varphi$. By \S.\,146, (3), if $\omega_1, \omega_2$ is a basis of $\mathfrak{o}$, we have to determine the integers $a_{1,1}, a_{1,2}, a_{2,2}$ in such a way that $a_{1,1}, a_{2,2}$ are positive and as small as possible, and that
\[ \alpha_1 = a_{1,1}\omega_1, \quad \alpha_2 = a_{1,2}\omega_1 + a_{2,2}\omega_2 \]
become divisible by $\varphi$. If we therefore set $(\omega_1, \omega_2) = (1, \Theta)$, then

\weq{1}{\alpha_1 = a_{1,1}, \quad \alpha_2 = a_{1,2} + a_{2,2}\Theta.}

Thus $a_{1,1}$ is the least positive rational number divisible by $\varphi$, and $a_{2,2}$ is the least positive number for which $a_{2,2}\Theta$ becomes congruent to a rational number according to the modulus $\varphi$.

We wish to apply this to the case where $\varphi$ is a prime functional $\pi$. Such a functional can here be only of the first or of the second degree, according as its absolute norm is equal to $p$ or $=p^2$. We therefore distinguish the two cases:

1) \quad $N_a(\pi) = p$

or

2) \quad $N_a(\pi) = p^2.$

Since here $N(p) = p^2$, in the first case $p$ is decomposable into two prime factors, while in the second case $p$ contains only one prime factor; that is, $p$ is still to be regarded as a prime in the field $\Omega$ itself.

If we now determine the basis (1) for $\varphi = \pi$, we obtain first of all $a_{1,1} = p$, and $a_{2,2}$ is the least positive integral rational number for which the product $a_{2,2}\Theta$ becomes congruent to a rational number modulo $\pi$. But for this, by \S.\,150, 3, the necessary and sufficient condition is
\[ (a_{2,2}\Theta)^p \equiv a_{2,2}\Theta \pmod{\pi}, \]
or, since by \gesperrt{Fermat}'s theorem $a_{2,2}^p \equiv a_{2,2}$,

\weq{2}{a_{2,2}(\Theta^p - \Theta) \equiv 0 \pmod{\pi}.}

Accordingly either $\Theta^p - \Theta \equiv 0$, and then $a_{2,2} = 1$, or $\Theta^p - \Theta$ is not $\equiv 0$, and then $a_{2,2} = p$.

In the first case $\Theta$, and with it every number $\omega$, is congruent to a rational number. There are therefore no more than $p$ numbers incongruent according to the modulus $\pi$, and consequently (\S.\,148, 3.) $p^f = p$, $f = 1$, that is, $p$ is decomposable into two prime factors of the first degree.
